\paragraph{Lee--Yang 分支图的 wreath 谱几何、dyadic affine Tate sheet 与 ZG 主极点残数}
在 Lee--Yang 有限分支图的超立方体地址、层间正则四叉结构与 Zeckendorf--Gödel 素数码的 Dirichlet 密度已经确定之后，还可把 H.158--H.164 的终端几何进一步压缩为四项可直接调用的后果。它们分别固定有限层 branch 图的全自同构、谱与 Ihara--Bass 数据，branch 塔逆极限的 affine Tate sheet 结构，以及 \(D_{\mathrm{ZG}}\) 在 \(s=1\) 处的主极点常数。

\begin{theorem}[Lee--Yang 第 \texorpdfstring{$n$}{n} 层分支图的超八面体花环对称]\label{thm:xi-time-part62dh-leyang-branch-hyperoctahedral-wreath}
沿用定理 \ref{thm:xi-time-part62d-leyang-branch-graph-gaussian-spectrum} 的记号。则有自然群同构
\[
\operatorname{Aut}(G_n)
\cong
\bigl((\ZZ/2\ZZ)^{2n}\rtimes S_{2n}\bigr)\wr S_5.
\]
\end{theorem}
\begin{proof}
这正是定理 \ref{thm:xi-time-part62d-leyang-branch-graph-automorphism-wreath} 的内容。证毕。
\end{proof}

\begin{corollary}[Lee--Yang 第 \texorpdfstring{$n$}{n} 层分支图的谱、Kirchhoff 森林与 Ihara--Bass 闭式]\label{cor:xi-time-part62dh-leyang-branch-spectrum-forest-ihara}
沿用定理 \ref{thm:xi-time-part62d-leyang-branch-graph-gaussian-spectrum} 的记号，记组合 Laplacian 为
\[
L_n:=2nI-A_n.
\]
则成立：
\begin{enumerate}
  \item 对每个 \(0\le k\le 2n\)，有
  \[
  \operatorname{mult}_{A_n}(2n-2k)=5\binom{2n}{k},
  \qquad
  \operatorname{mult}_{L_n}(2k)=5\binom{2n}{k}.
  \]
  \item 若记 \(\tau^\sharp(G_n)\) 为在每个连通分量内各取一棵生成树所得的组分保持生成森林数，则
  \[
  \tau^\sharp(G_n)
  =
  \Biggl(
  2^{2^{2n}-2n-1}
  \prod_{k=1}^{2n}k^{\binom{2n}{k}}
  \Biggr)^5.
  \]
  \item Ihara--Bass \(\zeta\) 函数满足
  \[
  \bigl(Z_{G_n}^{\mathrm{Ih}}(u)\bigr)^{-1}
  =
  \Biggl[
  (1-u^2)^{(2n-2)2^{2n-1}}
  \prod_{k=0}^{2n}
  \bigl(1-(2n-2k)u+(2n-1)u^2\bigr)^{\binom{2n}{k}}
  \Biggr]^5.
  \]
\end{enumerate}
\end{corollary}
\begin{proof}
第一条的邻接谱重数正是定理 \ref{thm:xi-time-part62d-leyang-branch-graph-gaussian-spectrum} 的第一条；再由
\[
L_n=2nI-A_n
\]
立即得到 Laplacian 谱重数公式。第二条即定理 \ref{thm:xi-time-part62d-leyang-branch-kirchhoff-forest}。第三条把定理 \ref{thm:xi-time-part62d-leyang-branch-ihara-zeta} 改写为逆函数形式便得。证毕。
\end{proof}

\begin{theorem}[Lee--Yang 分支塔逆极限的 affine Tate sheet 分解与层间正则四重覆盖]\label{thm:xi-time-part62dh-leyang-branch-affine-tate-regular-cover}
沿用推论 \ref{cor:xi-terminal-zm-leyang-finite-branch-regular-4ary-address} 的记号，并记
\[
\mathcal R_\infty^{(i)}:=\varprojlim_n\Pi_{i,n},
\qquad
\mathcal R_\infty:=\bigsqcup_{i=0}^{4}\mathcal R_\infty^{(i)}.
\]
则对每个 \(i\in\{0,\dots,4\}\) 与每个 \(n\ge 0\) 都有：
\begin{enumerate}
  \item 映射
  \[
  [2]:\Pi_{i,n+1}\longrightarrow \Pi_{i,n}
  \]
  是正则四重覆盖，并且
  \[
  \operatorname{Deck}\bigl([2]:\Pi_{i,n+1}\to \Pi_{i,n}\bigr)
  \cong
  E_{\min}[2]
  \cong
  (\ZZ/2\ZZ)^2.
  \]
  \item 一旦在第 \(i\) 个分支上选定相容提升 \(\mathbf P_i\in\mathcal R_\infty^{(i)}\)，便有
  \[
  \mathcal R_\infty^{(i)}\cong \mathbf P_i+T_2(E_{\min}),
  \qquad
  \mathcal R_\infty
  \cong
  \bigsqcup_{i=0}^{4}\bigl(\mathbf P_i+T_2(E_{\min})\bigr).
  \]
\end{enumerate}
特别地，Lee--Yang dyadic 分支塔的逆极限恰由五片彼此不交的 affine \(T_2(E_{\min})\)-sheet 组成。
\end{theorem}
\begin{proof}
推论 \ref{cor:xi-terminal-zm-leyang-finite-branch-regular-4ary-address} 的第一条已经给出
\[
[2]:\Pi_{i,n+1}\to\Pi_{i,n}
\]
为正则四重覆盖，且每个纤维在平移作用下都是 \(E_{\min}[2]\) 的主齐次空间。相应的 deck 变换正由核 \(E_{\min}[2]\) 的平移实现，因此 deck 群自然同构于 \(E_{\min}[2]\cong(\ZZ/2\ZZ)^2\)。这证明了第一条。

第二条则由定理 \ref{thm:xi-time-part62d-leyang-branch-inverse-limit-five-torsors} 与推论 \ref{cor:xi-terminal-zm-leyang-finite-branch-regular-4ary-address} 的第三条直接给出。证毕。
\end{proof}

\begin{corollary}[Zeckendorf--Gödel Dirichlet 级数主极点的自然密度残数]\label{cor:xi-time-part62dh-zg-dirichlet-density-residue}
\leanverified{paper\_xi\_time\_part62dh\_zg\_dirichlet\_density\_residue}
记
\[
D_{\mathrm{ZG}}(s):=\sum_{n\in\mathrm{ZG}}n^{-s}
\qquad (\Re(s)>1).
\]
则有
\[
\lim_{\sigma\downarrow 1}(\sigma-1)D_{\mathrm{ZG}}(\sigma)
=
\delta_{\mathrm{ZG}}=\delta(\mathrm{ZG}),
\qquad
\lim_{\sigma\downarrow 1}\frac{D_{\mathrm{ZG}}(\sigma)}{\zeta(\sigma)}
=
\delta_{\mathrm{ZG}}.
\]
因而 \(D_{\mathrm{ZG}}(s)\) 在 \(s=1\) 处具有留数 \(\delta(\mathrm{ZG})\) 的简单极点。
\end{corollary}
\begin{proof}
第一条正是定理 \ref{thm:xi-time-part62a-zg-simple-pole-density-residue}，第二条正是推论 \ref{cor:xi-time-part62a-zg-zeta-normalized-density-limit}。二者合并即得所述主极点残数公式。证毕。
\end{proof}

\noindent
由此，H.158--H.164 的 Lee--Yang 主线在结论层被再次压缩为一组彼此闭合的终端对象：有限层 branch 图的超八面体 wreath 对称、显式谱/森林/Ihara 数据与 dyadic 逆极限的五片 affine Tate sheet 彼此兼容，而 Zeckendorf--Gödel 像集的自然密度则同时外显为其 Dirichlet 级数在 \(s=1\) 的主极点残数。于是，超立方体地址、\(2\)-进仿射塔与 Euler 型解析常数便在同一审计框架内合流。

\endinput
